\documentclass[12pt,a4paper]{article}
\usepackage[utf8]{vietnam}
\usepackage{amsmath,amssymb}
\usepackage{tabularx}
\usepackage{tikz}
\usetikzlibrary{calc,angles,quotes,arrows.meta,patterns}
\usepackage{xcolor}
\usepackage{enumitem}
\usepackage{fancyhdr}
\usepackage[top=1.2cm, bottom=1.2cm, left=1.5cm, right=1.5cm]{geometry}
\usepackage[dethi]{ex_test}
\SAOPTN{widthSA=2.6,heightSA=0.62}

\setlength{\headheight}{15pt}
\setlength{\parskip}{0pt}

\pagestyle{fancy}
\fancyhf{}
\rhead{\textit{Môn Toán 6 -- Đề số 13}}
\lhead{\textit{Bộ đề ôn tập Giữa học kì I}}
\cfoot{\thepage}

\begin{document}

% --- KHUNG TIÊU ĐỀ ĐỀ THI ---
\noindent
\begin{tabularx}{\linewidth}{X r}
\textbf{TRƯỜNG THCS: \dotfill} & \textbf{ĐỀ KIỂM TRA ĐỊNH KỲ GIỮA HỌC KỲ I} \\
\textbf{Họ và tên: \dotfill} & \textbf{NĂM HỌC 2026 -- 2027} \\
\textbf{Lớp: \dotfill \quad SBD: \dotfill} & \textbf{Môn: TOÁN 6} \\
\hline
\end{tabularx}
\vspace{2mm}

\begin{center}
  \textbf{\large ĐỀ ÔN TẬP GIỮA HỌC KỲ I -- ĐỀ SỐ 13}\\
  \textit{Thời gian làm bài: 90 phút (Không kể thời gian phát đề)}
\end{center}
\vspace{1mm}

% --- PHẦN I ---
\noindent\textbf{\large PHẦN I. CÂU TRẮC NGHIỆM NHIỀU PHƯƠNG ÁN LỰA CHỌN (3,0 điểm)}

\textit{Thí sinh trả lời từ câu 1 đến câu 12. Mỗi câu hỏi thí sinh chỉ chọn một phương án duy nhất.}

\vspace{1mm}
\begin{ex}
Số vừa chia hết cho 2 vừa chia hết cho 3 là:
\choice
{$5019$}
{$5020$}
{$5021$}
{\True $5022$}
\end{ex}

\begin{ex}
Hình vuông $ABCD$ có cạnh $AB = 3\text{ cm}$. Diện tích của hình vuông đó bằng:
\choice
{$3\text{ cm}^2$}
{\True $9\text{ cm}^2$}
{$12\text{ cm}^2$}
{$81\text{ cm}^2$}
\end{ex}

\begin{ex}
Số nào sau đây không phải số nguyên tố?
\choice
{$11$}
{\True $33$}
{$17$}
{$31$}
\end{ex}

\begin{ex}
Ước chung lớn nhất của 13 và 4 là:
\choice
{\True $1$}
{$3$}
{$4$}
{$12$}
\end{ex}

\begin{ex}
Chọn câu khẳng định đúng nhất:
\choice
{Hình thang cân là hình thang có hai góc kề một đáy bằng nhau}
{Trong hình thang cân, hai cạnh bên bằng nhau}
{Trong hình thang cân, hai đường chéo bằng nhau}
{\True Cả A, B, C đều đúng}
\end{ex}

\begin{ex}
Tập hợp $\mathbb{N}$ tất cả các số tự nhiên là:
\choice
{$\{0; 1; 2\}$}
{$\{0; 1; 2; 3; 4\}$}
{\True $\{0; 1; 2; 3; \dots\}$}
{$\{0; 1; 2; 3; 4; 5; 6; 7; 8; 9\}$}
\end{ex}

\begin{ex}
Thứ tự thực hiện phép tính nào sau đây là đúng đối với biểu thức có dấu ngoặc?
\choice
{$[] \to () \to \{\}$}
{$[] \to \{\} \to ()$}
{$\{\} \to [] \to ()$}
{\True $() \to [] \to \{\}$}
\end{ex}

\begin{ex}
Điều kiện của số tự nhiên $x$ để biểu thức $A = 12 + 14 + 16 + x$ chia hết cho 2 là:
\choice
{\True $x$ là số tự nhiên chẵn}
{$x$ là số tự nhiên lẻ}
{$x$ là số tự nhiên bất kì}
{$x \in \{0; 2; 4; 6; 8\}$}
\end{ex}

\begin{ex}
Khẳng định nào sau đây là sai?
\choice
{Hình chữ nhật và hình thoi có các cặp cạnh đối song song với nhau}
{Hình lục giác đều có ba đường chéo chính bằng nhau}
{Hình bình hành có các cạnh đối bằng nhau}
{\True Hình thoi có hai đường chéo bằng nhau}
\end{ex}

\begin{ex}
Phát biểu nào sau đây là sai?\\Trong một hình bình hành:
\choice
{Hai đường chéo cắt nhau tại trung điểm mỗi đường}
{Các cạnh đối song song với nhau}
{\True Bốn góc ở các đỉnh bằng nhau}
{Các cạnh đối bằng nhau}
\end{ex}

\begin{ex}
Trong các số sau, số nào chia hết cho 2?
\choice
{$803$}
{$135$}
{$1975$}
{\True $358$}
\end{ex}

\begin{ex}
Giá một chiếc khẩu trang N95 là $30\,000$ đồng. Số tiền cần có để mua 20 chiếc khẩu trang N95 đó là:
\choice
{$30\,020$ đồng}
{$60\,000$ đồng}
{\True $600\,000$ đồng}
{$500\,000$ đồng}
\end{ex}


\vspace{1mm}
\noindent\textbf{BẢNG TRẢ LỜI TRẮC NGHIỆM PHẦN I:}
\begin{center}
\begin{tabular}{|c|c|c|c|c|c|c|c|c|c|c|c|c|}
\hline
\textbf{Câu} & \textbf{1} & \textbf{2} & \textbf{3} & \textbf{4} & \textbf{5} & \textbf{6} & \textbf{7} & \textbf{8} & \textbf{9} & \textbf{10} & \textbf{11} & \textbf{12} \\
\hline
\textbf{Chọn} & \rule{0pt}{0.65cm} & & & & & & & & & & & \\
\hline
\end{tabular}
\end{center}
\vspace{2mm}

\setcounter{ex}{0}
\noindent\textbf{\large PHẦN II. CÂU TRẮC NGHIỆM ĐÚNG SAI (2,0 điểm)}

\textit{Thí sinh trả lời câu 1 và câu 2. Trong mỗi ý a), b), c), d) ở mỗi câu, thí sinh chọn đúng hoặc sai.}

\vspace{1mm}
\begin{ex}
Cho hình vuông $ABCD$ có cạnh $AB = 30\text{ cm}$. Xét tính đúng/sai của các khẳng định sau:
\choiceTF
{\True Các góc của hình vuông bằng nhau và bằng $90^\circ$}
{Độ dài đường chéo $AC$ lớn hơn $BD$}
{Độ dài cạnh $BC = 20\text{ cm}$}
{Biết hình vuông $ABCD$ được ghép bởi 4 hình vuông nhỏ bằng nhau, cạnh của mỗi hình vuông nhỏ có độ dài là $15\text{ dm}$}
\end{ex}

\begin{ex}
Cho biểu thức $P = 2 \cdot 5^2 - 32 : 2^3$. Xét tính đúng/sai của các khẳng định sau:
\choiceTF
{Giá trị của lũy thừa $5^2$ bằng 10}
{\True Biểu thức $P$ có giá trị bằng 46}
{Giá trị của biểu thức $P$ chia hết cho cả 2 và 5}
{\True Giá trị của biểu thức $P$ là một hợp số}
\end{ex}


\noindent\textbf{BẢNG TRẢ LỜI PHẦN II:}
\begin{center}
\begin{tabular}{|c|c|c|c|c||c|c|c|c|c|}
\hline
\textbf{Câu 1} & \textbf{a)} & \textbf{b)} & \textbf{c)} & \textbf{d)} & \textbf{Câu 2} & \textbf{a)} & \textbf{b)} & \textbf{c)} & \textbf{d)} \\
\hline
\textbf{Đ/S} & \rule{0pt}{0.6cm} & & & & \textbf{Đ/S} & & & & \\
\hline
\end{tabular}
\end{center}
\vspace{2mm}

\setcounter{ex}{0}
\noindent\textbf{\large PHẦN III. CÂU HỎI TRẢ LỜI NGẮN (2,0 điểm)}

\textit{Thí sinh trả lời câu 1 và câu 2. Điền kết quả số vào ô trống tương ứng.}

\vspace{1mm}
\begin{ex}
Tìm số tự nhiên $x$ lớn nhất thỏa mãn $48 \ \vdots \ x$ và $72 \ \vdots \ x$. \shortans{24}
\end{ex}

\begin{ex}
Một hình thoi có diện tích bằng $30\text{ cm}^2$ và độ dài một đường chéo bằng $6\text{ cm}$. Tính độ dài đường chéo còn lại của hình thoi (theo đơn vị $\text{cm}$). \shortans{10}
\end{ex}


\setcounter{ex}{0}
\noindent\textbf{\large PHẦN IV. TỰ LUẬN (3,0 điểm)}

\textit{Thí sinh trình bày chi tiết lời giải vào giấy làm bài.}

\vspace{1mm}
\begin{ex} (1,0 điểm)
Thực hiện các phép tính sau (Tính nhanh nếu có thể):
\begin{enumerate}[label=\textbf{\alph*)}, leftmargin=0.8cm, itemsep=2pt]
\item (0,25 điểm) $15 + 187 + 85$
\item (0,25 điểm) $2024 - [38 - (4 - 1)^3]$
\item (0,25 điểm) $27 \cdot 39 + 27 \cdot 63 - 2 \cdot 27$
\item (0,25 điểm) $210 : \{38 - [12 + 3 \cdot (32 : 2^3)]\}$
\end{enumerate}
\end{ex}
\vspace{1.5mm}

\begin{ex} (1,0 điểm)
Tìm số tự nhiên $x$, biết:
\begin{enumerate}[label=\textbf{\alph*)}, leftmargin=0.8cm, itemsep=2pt]
\item (0,25 điểm) $19 - x = 5$
\item (0,25 điểm) $2 \cdot 3^x - 5 = 7^2$
\item (0,25 điểm) $81 - (3x + 5) = 4^3$
\item (0,25 điểm) $(19x + 2 \cdot 5^2) : 14 = (13 - 8)^2 - 4^2$
\end{enumerate}
\end{ex}
\vspace{1.5mm}

\begin{ex} (0,5 điểm)
Nhà bạn Nam có mảnh vườn hình chữ nhật, chiều dài $12\text{ m}$, chiều rộng $8\text{ m}$. Ở giữa vườn có khu nghỉ ngơi hình chữ nhật chiều dài $6\text{ m}$, chiều rộng $2\text{ m}$. Phần còn lại của khu vườn dùng để trồng cây.
\begin{center}
\begin{tikzpicture}[scale=0.5, font=\small]
  \draw[thick, fill=green!15] (0,0) rectangle (12,8);
  \draw[thick, fill=orange!20] (3,3) rectangle (9,5);
  \node at (6,4) {\footnotesize Khu nghỉ ngơi};
  \node[above] at (6,8) {$12\text{ m}$};
  \node[right] at (12,4) {$8\text{ m}$};
  \node[below] at (6,3) {$6\text{ m}$};
  \node[left] at (3,4) {$2\text{ m}$};
\end{tikzpicture}
\end{center}
\begin{enumerate}[label=\textbf{\alph*)}, leftmargin=0.8cm, itemsep=2pt]
\item (0,25 điểm) Tính diện tích toàn bộ mảnh vườn?
\item (0,25 điểm) Biết mỗi mét vuông cây trồng hết $15\,000$ đồng để chăm bón hàng tháng, hỏi bạn Nam cần bao nhiêu tiền chăm bón hàng tháng?
\end{enumerate}
\end{ex}
\vspace{1.5mm}

\begin{ex} (0,5 điểm)
Bài toán nâng cao về lũy thừa và chia hết:
\begin{enumerate}[label=\textbf{\alph*)}, leftmargin=0.8cm, itemsep=2pt]
\item (0,25 điểm) Cho số tự nhiên $x$ thỏa mãn $3^{x+1} + 3^{x+2} - 2 \cdot 3^x = 270$. Hỏi $x$ là số nguyên tố hay hợp số?
\item (0,25 điểm) Tìm hai số tự nhiên biết tích của chúng bằng 144 và ƯCLN của chúng bằng 6.
\end{enumerate}
\end{ex}
\vspace{1.5mm}


\begin{center}
\rule{4cm}{0.4pt}\\
\textit{--- HẾT ---}
\end{center}
\end{document}
